\chapter{Project Management} \label{chapter6}
\chaptermark{Project Management}
\section{Gantt Chart and Project Planning}
\begin{figure}[h!]
\begin{ganttchart}{1}{12}
  \gantttitle{First semester}{12} \\
  \gantttitlelist{1,...,12}{1} \\
  \ganttbar{Research \& Preparation}{1}{12} \\
  \ganttbar{Prototype Development}{11}{12} \\
  \ganttbar{Evaluation}{8}{10} \\
  \ganttbar{Writing \& Finalisation}{1}{12}
\end{ganttchart}
\caption{Gantt chart outlining work plan for the first semester.}
\label{fig:gantt-first-semester}
\end{figure}

\begin{figure}[h!]
\centering
\begin{ganttchart}{1}{12}
  \gantttitle{Second semester semester}{12} \\
  \gantttitlelist{1,...,12}{1} \\
  \ganttbar{Preparation}{1}{2} \\
  \ganttbar{Prototype Development}{1}{3} \\
  \ganttbar{Design \& Implementation}{2}{12}\\
  \ganttbar{Writing \& Finalisation}{1}{12}
\end{ganttchart}
\caption{Gantt chart outlining work plan for the second semester.}
\label{fig:gantt-second-semester}
\end{figure}
The two Gantt charts represent the workflow of both semesters. The first semester focused on Writing and research. The second semester focused on design and implementation.


\begin{table}[h!]
\centering
\begin{tabular}{|p{4cm}|p{3cm}|p{7cm}|}
\hline
\textbf{Risk} & \textbf{Likelihood / Impact} & \textbf{Mitigation Strategy} \\
\hline
Model fails to converge or underperforms & Medium / High & Use pre-trained embeddings as fallback; incremental tuning; monitor loss curves. \\
\hline
Insufficient dataset diversity & Medium / Medium & Combine FMA and GTZAN; create synthetic augmentations to simulate real-world transformations. \\
\hline
Computational constraints & Low / Medium & Optimise batch sizes, use Google Colab Pro or university GPU resources. \\
\hline
Time overruns & Medium / High & Maintain weekly milestones; use iterative development; prioritise feature completeness over optional enhancements. \\
\hline
Incorrect assumptions about robustness & Medium / Medium & Conduct early feasibility tests; refine methodology based on preliminary findings. \\
\hline
Augmentations produce unrealistic or harmful distortions 
& Medium / Medium 
& Validate augmentation strengths against real short-form audio; tune augmentation parameters to avoid training on unrealistic extremes. \\
\hline
FAISS indexing fails to scale or gives unstable retrieval results 
& Low / Medium 
& Test multiple index types (Flat, IVFFlat, HNSW); monitor recall at k; fall back to brute-force cosine search for smaller experiments. \\
\hline
Spectrogram extraction or preprocessing pipeline breaks 
& Medium / Low 
& Validate audio loading and spectrogram generation on smaller subsets; use robust libraries (Librosa) with version control to maintain reproducibility. \\
\hline
Embedding dimensionality mismatch between models 
& Medium / Medium 
& Standardise embedding shapes or use projection layers; maintain strict documentation for embedding configurations. \\
\hline
Overfitting in contrastive model 
& Medium / High 
& Use large augmentation variety, early stopping, dropout, and validation subsets; monitor representation collapse via embedding variance metrics. \\
\hline
Versioning or dependency conflicts 
& Medium / Medium 
& Use a pinned environment file (requirements.txt); track library versions; test the pipeline on a clean environment before final submission. \\
\hline
File corruption or dataset download failure 
& Low / High 
& Store mirrored copies on cloud storage; add checksum verification; maintain dataset-loading fallback routines. \\
\hline
Difficulty interpreting evaluation results or insufficient evaluation coverage 
& Medium / Medium 
& Use well-defined metrics (precision, recall, F1, robustness curves); prepare multiple test subsets with varying transformation strengths. \\
\hline
Risk of producing misleading or overstated results 
& Low / High 
& Clearly document limitations; avoid claims about real-world performance beyond tested scenarios; follow ethical reporting practices. \\
\hline
\end{tabular}
\caption{Risk analysis with associated mitigation strategies.}
\end{table}

\section{Professional, Legal, Ethical, and Social Considerations (PLES)}
This project engages several important PLES topics:


\subsection*{Professional}
The work follows professional standards in computing research, including 
documentation quality, reproducibility, and adherence to the BCS Code of Conduct. 
Open datasets and transparent evaluation protocols ensure responsible research practice.

\subsection*{Legal}
The project processes only openly licensed audio (FMA, GTZAN). No copyrighted 
commercial music is used, preventing violations of copyright law. Data processing 
contains no personal or identifiable information, maintaining GDPR compliance.

\subsection*{Ethical}
Use of synthetic transformations avoids ethically problematic data collection.  
Automated copyright detection systems have broader ethical implications, such as 
false positives, platform moderation fairness, and potential misuse. These concerns 
are considered and discussed in the concluding sections.

\subsection*{Social}
Short-form media ecosystems rely heavily on creator freedoms and remix culture.  
Improved detection systems may impact creators, rights holders, and digital 
platforms. WaveID aims to support fair remuneration while avoiding unnecessary 
restrictions on creative expression.

\section{GenAI Summary}
GenAI will not be used in the project as any datasets, preprocessing, embedding, and any other operations will be handled independently through Python, PyTorch, Librosa and FAISS. It is also not necessary in development of the Deliverable.

\section{Risk Analysis and Mitigation}
The following risks were identified during the planning stage:
